\section{Background and Motivation}
\label{sec:background}

\subsection{Expert parallelism in MoE}
\label{sec:epbackground}
A Mixture-of-Experts layer routes each input token to a subset of expert networks.
Under expert parallelism (EP), the experts are partitioned across GPU ranks. Dispatch
transfers token activations to the selected expert ranks, the experts compute their
outputs, and Combine returns and merges these contributions at the token
owner~\cite{deepep,ncclep}. Each token selects up to $\topk$ experts, so a single
input may require contributions from several ranks.

Two properties of this pattern matter for what follows. First, routing is
input-dependent: the number of records a rank sends to a particular destination is not
known until the gate has run, so communication carries metadata describing valid
counts, destinations, and output placement alongside the payload. Second, the
communication is genuinely many-to-many and it straddles a hardware boundary. Within a
node, ranks exchange data over a scale-up interconnect such as NVLink; across nodes,
they use a scale-out transport such as RDMA. Communication libraries manage both paths
together with the staging buffers that prepare data for expert computation.

\subsection{Three dependencies, two synchronization points}
\label{sec:epsynchronization}
Correctness here rests on three distinct dependencies, and conflating them is the
source of most confusion in this area. Before a sender writes, the destination's
communication resources must be ready to receive. Before a consumer reads, the
required data must be visible to it. Before a buffer is reused, its previous consumers
must have finished reading it. These are not one condition observed at three times: a
transport operation's local progress, the arrival of its data at the destination, and
the end of its consumers' reads are separate events, and a given barrier or completion
signal enforces some combination of them rather than all three.

DeepEP V2's cross-node Direct paths make both transfer boundaries
explicit~\cite{deepepcode}. An entry barrier coordinates the participating ranks
before payload injection. Dispatch then overlaps count exchange and prefix computation
with data movement, and Combine reuses the routing Dispatch established. After the
data-transfer loop, a separate tail stage waits for outstanding asynchronous memory
operations, flushes the GIN queues, joins the local workers, and only then exchanges
completion signals across the EP group; those signals release the subsequent copy or
reduction epilogue. Because staging provides known landing addresses, an unknown final
compact layout is not by itself the reason for the entry barrier.

This organization is specific to the implementation, and we are careful not to
generalize it. Direct handles cross-node communication, including local NVLink access
where available, while Hybrid performs hierarchical scale-out transfer with scale-up
forwarding and ends Dispatch with a scale-up-only barrier. NCCL EP HT instead performs
a routing AllGather and metadata scan before Dispatch, and LL can begin data transfer
with no independent readiness barrier at
all~\cite{deepepcode,ncclhtcode,ncclllcode}. Our subject is cross-node Direct
Dispatch and Combine, the path that exposes both boundaries; Table~\ref{tab:paths}
records where each other path falls.

\subsection{Why small batches change the accounting}
\label{sec:smallbatch}
Payload transmission time falls with batch size. Readiness announcements, local joins,
and completion notifications do not: their cost is set by traversal and endpoint
processing, not by payload bytes. Fixed coordination therefore grows as a fraction of
a shrinking operator, and decode sits squarely in that regime, because each active
request contributes roughly one token per step spread across data-parallel ranks.
Since $T/\text{rank} \approx \text{concurrency}/\text{EP}$, raising concurrency does
not leave the region: it is a property of the serving configuration rather than a
tuning oversight.

Our H200 measurements (\secref{sec:evidence}) quantify both boundaries on the Direct
path itself. At EP8 and eight tokens per rank, entry synchronization takes about
15.1\us{} and the completion-signal window 13.5--14.3\us{}, amounting to roughly
20\% and 18\% of the corresponding traced operator. Sweeping to 32, 64, and 128 tokens
per rank leaves the absolute durations essentially where they were---entry stays within
14.8--15.4\us{}---while their shares fall monotonically to about 5\%. The cost did not
go away; the denominator grew. This frames the
question the rest of the paper addresses: must readiness and completion coordination
remain sequential with payload transfer, or can they overlap it while preserving data
visibility and buffer-lifetime guarantees?
